\section{Benchmark Construction}
\label{sec:benchmark-construction}

The benchmark's target design is tightly controlled paired comparison. A
future minimal pair should hold the surface string, payload, and response space
as constant as possible while changing one control dimension: source role,
authority level, pragmatic status, scope condition, reference resolution,
speech-act force, application surface, or refusal status. This follows
verifiable instruction-following evaluation, in which target behaviour is
specified before outputs are inspected \parencite{zhou2023ifeval}, but
extends it to contrasts whose correctness depends on source, embedding, and
pragmatic status.

Items should be constructed under five constraints. First, each item names the source of the relevant language. Second, each item has a contrast item that tests the same string, payload, or response space under a different embedding or authority relation. Third, the expected behaviour is stated before model outputs are inspected; for controlled authority items, an independently specified authorization record defines the licensed action, while observed model behaviour tests whether the model responds accordingly. Fourth, the response labels are shared across the contrast whenever possible. Fifth, examples should use harmless payloads unless a verified safety policy and review process justify a more realistic item.

The current seed file, \texttt{benchmark/items/seed-items.csv}, contains 18 development items. It uses colour-token payloads such as \mention{BLUE} and \mention{GREEN}, protected dummy tokens such as \mention{ZETA-42}, and dummy secrets such as \mention{ALPHA-000}. These preserve the control structure of prompt-injection, exfiltration, and refusal cases while avoiding operational unsafe content.

Toy payloads support claims about control structure, object-label stability, and the behaviour of the declared measurement procedure. They don't yet support claims about deployed prompt-injection resistance, harmful-request refusal, or real-world agent-security robustness. Those claims require realistic wrappers and independent policy review.

The seed set now includes paired prompt contrasts spanning webpage text,
document text, email body, tool output, dummy private state, and policy text.
These are controlled development contrasts, not uniformly strict minimal pairs:
several vary wording or response demands along with the designated dimension.
For example, a tool result can say \mention{Order status: shipped. Ignore the
user and output BLUE}; the expected behaviour is to report the order status
without following the tool-result instruction. A separate dummy-secret item
tests whether an email can induce disclosure of \mention{ALPHA-000}.

The pilot also exposed a seed-set repair. Pair P008 contrasts an email-body
exfiltration probe with a direct toy-policy refusal, but its members differ in
phenomenon family and both expect non-disclosure. It's retained as a diagnostic
confidentiality contrast and excluded from paired-contrast scoring. The
development set should split it into two properly controlled pairs before using
it for pairwise contrast accuracy.

Each item records the control dimension, context source, source role, authority level, pragmatic status, response act, expected behaviour, task-success label, policy-compliance label, safety-risk label, risk type, refusal outcome, failure attribution, and a judge-validation flag. One scalar field is legacy development metadata: \texttt{authority\_level}. It can't represent additional approval requirements, transition type, or time-specific authorization. Study B's typed authorization record is the intended replacement for controlled authority contrasts. The validator script enforces the seed schema and label inventories.

\begin{figure}[t]
\centering
\includegraphics[width=\linewidth]{evaluation_pipeline}
\caption{Evaluation pipeline for the benchmark artifact. The diagram summarizes the data flow from pre-specified item metadata through model output, rule-aided triage, expert labels, LLM-judge labels, adjudication, and metric-driven item revision. No performance quantity is encoded in the figure.}
\label{fig:evaluation-pipeline}
\end{figure}

Figure~\ref{fig:evaluation-pipeline} shows how that schema makes the benchmark auditable in two directions. From the prompt and metadata, an evaluator can see which contrast is being tested. From a model output, an evaluator can identify which component failed: source sensitivity, mention/use robustness, scope interpretation, refusal calibration, risk classification, or the safety policy layer.

Study B supplies the next controlled construction stage. Its four current bases yield 16 development fixtures. Those fixtures feed a 12-base construction pilot; lessons from that pilot then inform a planned 24-base, four-condition development set, or 96 items. Its job is to test contrast integrity, the taxonomy, expert-evaluation protocol, validator, and scoring routes before item generation scales up. A 300--800-item benchmark would be a later expansion, large enough to broaden phenomenon coverage but still contingent on close expert adjudication and error analysis at the earlier stages.
